\documentclass[10pt,a4paper]{dinbrief}
\usepackage[latin1]{inputenc}
\usepackage{ngerman}
\usepackage{a4wide}
\usepackage{eurosym}
\usepackage{fancyhdr}
\usepackage{pdfpages}
\address{% 
	Dr. Tobias Griebe\\
	Gartenstra{\ss}e 29\\
	15366 Neuenhagen 
}
\backaddress{Dr. T. Griebe, Gartenstr. 29, 15366 Neuenhagen}
\signature{Dr. Tobias Griebe}
\place{Neuenhagen}
\nowindowrules \nobackaddressrule \nowindowtics \normaladdress
\date{\today}
\pagestyle{empty}
\begin{document} 
\begin{letter}  
{%
	ic! berlin brillen GmbH\\
	Wolfener Stra"se 32 F\\
	12681 Berlin\\
	Deutschland
}
\subject{\bf Bewerbung Head of IT}
\opening{Sehr geehrte Damen und Herren,}
auf der Suche einer beruflichen Ver"anderung bin ich mit gro"sem Interesse auf Ihre Stellenausschreibung Head of IT am Standort Berlin gesto"sen. 

In meinen bisherigen T"atigkeiten als Competence Center Leiter bei der adesso AG im Bereich Digitalisierung sowie zuvor als wissenschaftlicher Mitarbeiter an der Universit"at Duisburg-Essen konnte ich meine Kompetenzen bei Konzeption und Entwicklung von digitalen L"osungen im gesamten Spektrum Ideation, Anforderungsermittlung in fr"uhen Projektphasen bis hin zur Implementierung, Qualit"assicherung und Inbetriebnahme unter Beweis stellen.

Als Leiter der Arbeitsgruppe f"ur mobile Anwendungen und Prozesse am wirtschaftsnahen Lehrstuhl f"ur Software Engineering, insb. f"ur mobile Anwendungen der Universit"at Duisburg-Essen hatte ich Gelegenheit meine Kenntnisse in den Bereichen Requirements Engineering, Softwareentwicklung und Qualit"atsmanagement in Forschung und Lehre sowie in zahlreichen Industrieprojekten in der Praxis einzusetzen. Die Leitung von IT-Projekten und die Verwendung innovativer Technologien geh"orten dabei ebenso zum Projektalltag wie der Einsatz agiler Arbeitsmethoden und der Aufbau und die F"uhrung von Projektteams.

Meine T"atigkeit an der Universit"at m"undete im erfolgreichen Abschluss meiner Doktor\-arbeit zum Thema Methode und Technologie zur modellbasierten Automatisierung von Tests kontextsensitiver mobiler Anwendungen, in welcher ich den vollst"andigen Workflow vom modellbasierten Testfallentwurf auf Basis von UML-Aktivit"atsdiagrammen bis hin zur vollautomatischen Testausf"uhrung f"ur mobile Anwendungen auf der Plattform Android konzipiert, implementiert und evaluiert habe.

\clearpage
Den Aufbau und die fachliche und disziplinarische F"uhrung eines Teams von Beratern im Themenfeld Digitalisierung und digitale Transformation konnte ich bei der adesso AG als Competence Center Leiter fortsetzen. Mein Hauptverantwortungsbereich bei der adesso AG ist die fachliche und disziplinarische F"uhrung von Mitarbeitern "uber den gesamten Lebenszyklus von Recruiting, Qualifikation, Laufbahnstufenentwicklung und Retention, der Auf- und Ausbau technischer und methodischer Kompetenzen und den Transfer auf die Mitarbeiter, Entwicklung von Mitarbeitern, identifizieren von Leistung und Potenzial eigener Mitarbeiter, ermitteln von  Entwicklungsbedarf, zielgerichteter Einsatz passender Personalentwicklungsinstrumente, Risikomanagement und vorausschauendes Eskalationsmanagement.


Neben dem Aufbau meines Teams befasse ich mich mit der Durchf"uhrung von Projekten im gesamten Handlungsfeld von Ideation, Identifikation digitalisierbarer Gesch"afts\-prozesse, Analyse von Digitalisierungspotenzial existierender und neuer Fachanforderungen, Requirements Engineering, der Erstellung von Grob- und Feinspezifikationen, der Koordination, Durchf"uhrung und Begleitung von IT-Projekten in strategischen und operativen Rollen, Erstellung von Handlungsempfehlungen bei Beschaffung, Implementierung und Inbetriebnahme technischer L"osungen sowie mit der Steuerung interner und externer Entwicklungsdienstleister. 

Bei dieser T"atigkeit bediene ich mich aus einem umfangreichen Methodenbaukasten zur Identifikation und Konkretisierung von Digitalisierungspotenzial (u.~a. Business Model Canvas, Stakeholder-Analyse, Touchpoint-Analyse, Customer Journey Mapping und Customer Safari, Interaction Room) sowie g"angiger Modellierungssprachen und Werkzeuge zur Dokumentation und Analyse von Prozessen (z.~B. BPMN, UML).

Mit Fokus auf die Umsetzung von (IT-)Projekten geh"oren der sichere Umgang mit klassischen und agilen Vorgehensmodellen (insb. Scrum) ebenso zu meinen Kompetenzen wie der Umgang mit State-of-the-Art Werkzeugen zur aktiven und passiven Begleitung von IT-Projekten (z.~B. Jira, Confluence oder MS-TFS).

Nun bin ich auf der Suche nach einem neuen Engagement in einem innovativen und dynamischen Umfeld und sehe den Herausforderungen des Head of IT mit Spannung entgegen.

Vor diesem Hintergrund bin ich davon "uberzeugt, durch meinen ausgepr"agten technischen und akademischen Hintergrund, meine eigenst"andige und qualit"atsorientierte Arbeitsweise sowie meine F"ahigkeit Mitarbeiter anzuleiten, zu entwickeln und zu motivieren zum Erfolg Ihres Unternehmens beizutragen.

Ich bin fr"uhestens ab April 2019 verf"ugbar. Aufgrund meiner Qualifikation meiner bisherigen Berufsfahrungen liegen meine Gehaltsvorstellungen bei 105.000 Euro/Jahr. 

"Uber eine Einladung zu einem pers"onlichen Gespr"ach freue ich mich sehr.

\closing{Mit freundlichen Gr"u"sen,}
\encl{Lebenslauf, Promotionsurkunde, Diplomurkunde, Zertifikat Licence to Lead F"uhrungskr"afteausbildung, Zertifikat Professional Scrum Product Owner (PSPO I), Zertifikat ISTQB Certified Tester Foundation Level}
\end{letter}
\end{document}